% ===============================================================
%  جعبه‌ابزار ریاضی ۲
% ===============================================================
\chapter{انتگرال دوگانه، تغییر متغیر و ژاکوبین}

از جلسه ۱۳ به بعد با \textbf{دو یا چند} متغیر تصادفی سروکار داریم: چگالی مشترک، چگالی حاشیه‌ای، کوواریانس، توزیع جمع و نسبت، آماره‌های ترتیبی و توزیع‌های مشترکاً گاوسی. ابزار اصلی همه این‌ها \textbf{انتگرال دوگانه} است. در این بخش، انتگرال دوگانه را از پایه مرور می‌کنیم و مهم‌ترین مهارت آن را تمرین می‌کنیم: \textbf{تعیین درست حدود انتگرال برای یک ناحیه}. سپس مختصات قطبی (و اثبات انتگرال گاوسی)، تغییر متغیر کلی با \textbf{ژاکوبین} و کاربردش در یافتن چگالی تبدیل‌های دومتغیره را می‌بینیم.

\section{انتگرال دوگانه: احتمال به عنوان حجم}
اگر $(X,Y)$ چگالی مشترک $f_{X,Y}(x,y)\ge0$ داشته باشند، احتمال اینکه نقطه تصادفی $(X,Y)$ در ناحیه $A$ از صفحه بیفتد، \textbf{حجم} زیر رویه $z=f_{X,Y}(x,y)$ و بالای ناحیه $A$ است:
\[
\Prob\big((X,Y)\in A\big)=\iint_Af_{X,Y}(x,y)\dd x\dd y,\qquad \iint_{\R^2}f_{X,Y}(x,y)\dd x\dd y=1 .
\]
شهود دقیقاً مانند حالت یک‌بعدی است: صفحه را به مستطیل‌های کوچک $\Delta x\times\Delta y$ تقسیم می‌کنیم؛ احتمال افتادن در هر مستطیل تقریباً $f(x_i,y_j)\Delta x\Delta y$ (ارتفاع × مساحت قاعده) است و انتگرال حد مجموع آن‌هاست.

\begin{center}
\renewcommand{\arraystretch}{1.6}
\begin{tabular}{c|c}
مفهوم (جلسه‌های ۱۳ تا ۱۷) & انتگرال دوگانه\\\hline
چگالی حاشیه‌ای & $f_X(x)=\int_{-\infty}^{\infty}f_{X,Y}(x,y)\dd y$\\
امید یک تابع & $\E g(X,Y)=\iint g(x,y)f_{X,Y}(x,y)\dd x\dd y$\\
کوواریانس & $\Cov(X,Y)=\E XY-\E X\,\E Y$\\
مقایسه دو متغیر & $\Prob(X<Y)=\iint_{x<y}f_{X,Y}\dd x\dd y$\\
توزیع جمع & $F_{X+Y}(z)=\iint_{x+y\le z}f_{X,Y}\dd x\dd y$
\end{tabular}
\end{center}

\section{انتگرال تکراری و قضیه فوبینی}
\begin{theorem}[فوبینی]
اگر $f$ روی مستطیل $[a,b]\times[c,d]$ انتگرال‌پذیر باشد (برای چگالی‌ها، که نامنفی‌اند، همواره برقرار است)، آنگاه
\[
\iint_{[a,b]\times[c,d]}f(x,y)\dd x\dd y=\int_a^b\Big(\int_c^df(x,y)\dd y\Big)\dd x=\int_c^d\Big(\int_a^bf(x,y)\dd x\Big)\dd y .
\]
\end{theorem}
یعنی انتگرال دوگانه را می‌توان به صورت \textbf{دو انتگرال یگانه پشت سر هم} حساب کرد. در انتگرال درونی، متغیر بیرونی مثل یک \textbf{عدد ثابت} رفتار می‌کند. شهود: حجم را به برش‌های نازک موازی محور $y$ تقسیم می‌کنیم؛ انتگرال درونی مساحت هر برش (به ازای $x$ ثابت) است و انتگرال بیرونی این مساحت‌ها را جمع می‌زند.

\textbf{حالت جداپذیر.} اگر $f(x,y)=g(x)h(y)$ و ناحیه مستطیل باشد، انتگرال به حاصل‌ضرب دو انتگرال یگانه تبدیل می‌شود: $\iint g(x)h(y)=\Big(\int g\Big)\Big(\int h\Big)$. در احتمال، این همان \textbf{استقلال} است: $X$ و $Y$ مستقل‌اند اگر و تنها اگر چگالی مشترک به صورت حاصل‌ضرب تابعی از $x$ و تابعی از $y$ تجزیه شود (روی یک ناحیه مستطیلی)؛ و آنگاه $\E[g(X)h(Y)]=\E g(X)\,\E h(Y)$.

\begin{example}[چگالی مشترک روی مربع واحد]
$f(x,y)=c(x+y)$ برای $0\le x,y\le1$. (الف) $c$، (ب) $f_X$، (ج) $\E X$، $\E XY$ و $\Cov(X,Y)$ را بیابید.
\begin{solution}
(الف) انتگرال درونی نسبت به $y$ (با $x$ ثابت): $\int_0^1(x+y)\dd y=x+\frac12$. سپس $\int_0^1\big(x+\frac12\big)\dd x=\frac12+\frac12=1$؛ پس $c=1$.

(ب) همان انتگرال درونی: $f_X(x)=\int_0^1(x+y)\dd y=x+\frac12$ برای $0\le x\le1$ (و به دلیل تقارن، $f_Y(y)=y+\frac12$).

(ج)
\[
\E X=\int_0^1x\Big(x+\frac12\Big)\dd x=\frac13+\frac14=\frac{7}{12},
\]
\[
\E XY=\int_0^1\!\!\int_0^1xy(x+y)\dd y\dd x=\int_0^1\Big(\frac{x^2}{2}+\frac{x}{3}\Big)\dd x=\frac16+\frac16=\frac13 ,
\]
(انتگرال درونی: $\int_0^1(x^2y+xy^2)\dd y=\frac{x^2}{2}+\frac x3$.)
\[
\Cov(X,Y)=\frac13-\Big(\frac{7}{12}\Big)^2=\frac{48-49}{144}=-\frac{1}{144}.
\]
کوواریانس منفی است و $X,Y$ مستقل نیستند: $x+y$ را نمی‌توان به صورت $g(x)h(y)$ نوشت.
\end{solution}
\end{example}

\section{ناحیه‌های غیرمستطیلی: تعیین حدود انتگرال}
در احتمال، ناحیه‌ها تقریباً هیچ‌گاه مستطیل نیستند: «$X<Y$»، «$X+Y\le z$»، «$0<y<x<1$». مهم‌ترین مهارت، نوشتن درست حدود است.

\textbf{دستورالعمل گام‌به‌گام:}
\begin{enumerate}
\item ناحیه را \textbf{رسم کنید}: همه نامساوی‌ها (از جمله تکیه‌گاه چگالی و شرط رویداد) را روی صفحه بکشید و اشتراکشان را سایه بزنید.
\item ترتیب را انتخاب کنید. اگر انتگرال درونی نسبت به $y$ است (\textbf{برش عمودی}): یک $x$ ثابت در نظر بگیرید و ببینید خط عمودی در آن $x$ از کجا وارد ناحیه می‌شود ($y=g_1(x)$) و از کجا خارج ($y=g_2(x)$).
\item حدود بیرونی، کوچک‌ترین و بزرگ‌ترین $x$ در کل ناحیه است (اعداد ثابت).
\item اگر مرز ناحیه در طول بازه $x$ «عوض شود» (مثلاً شکسته باشد)، انتگرال بیرونی را به چند قسمت بشکنید.
\end{enumerate}
\[
\iint_Af=\int_{x_{\min}}^{x_{\max}}\Big(\int_{g_1(x)}^{g_2(x)}f(x,y)\dd y\Big)\dd x
\qquad\text{یا}\qquad
\int_{y_{\min}}^{y_{\max}}\Big(\int_{h_1(y)}^{h_2(y)}f(x,y)\dd x\Big)\dd y .
\]
\textbf{قانون طلایی:} حدود انتگرال درونی می‌توانند به متغیر بیرونی وابسته باشند، اما حدود انتگرال بیرونی \textbf{همیشه عدد ثابت‌اند}.

\begin{example}[ناحیه مثلثی: بیشینه و کمینه دو نمونه یکنواخت]
$f(x,y)=2$ روی مثلث $0<y<x<1$ (این چگالی مشترک بیشینه و کمینه دو نمونه مستقل $\mathcal U(0,1)$ است؛ جلسه ۱۷). (الف) نرمال بودن را بررسی کنید. (ب) چگالی‌های حاشیه‌ای را بیابید. (ج) $\Prob(X>2Y)$ را حساب کنید.
\begin{solution}
(الف) برش عمودی: به ازای $x$ ثابت در $(0,1)$، $y$ از $0$ تا $x$ تغییر می‌کند:
\[
\int_0^1\!\!\int_0^x2\dd y\dd x=\int_0^12x\dd x=1 .
\]
(ب) $f_X(x)=\int_0^x2\dd y=2x$ برای $0<x<1$ (چگالی بیشینه). برای $f_Y$ برش افقی: به ازای $y$ ثابت، $x$ از $y$ تا $1$: $f_Y(y)=\int_y^12\dd x=2(1-y)$ (چگالی کمینه).

(ج) شرط $x>2y$ یعنی $y<\frac x2$. اشتراک با ناحیه: به ازای $x$ ثابت، $y$ از $0$ تا $\frac x2$ (چون $\frac x2<x$):
\[
\Prob(X>2Y)=\int_0^1\!\!\int_0^{x/2}2\dd y\dd x=\int_0^1x\dd x=\frac12 .
\]
\end{solution}
\end{example}

\begin{center}
\begin{tikzpicture}
\end{tikzpicture}\\
{\small ناحیه $0<y<x<1$ و یک برش عمودی: به ازای $x$ ثابت، $y$ از $0$ (خط پایین) تا $x$ (قطر) تغییر می‌کند. خط‌چین مرز $y=x/2$ است که رویداد $X>2Y$ را جدا می‌کند.}
\end{center}

\begin{example}[احتمال $X+Y\le1$]
برای چگالی $f(x,y)=x+y$ روی مربع واحد (مثال بخش فوبینی)، $\Prob(X+Y\le1)$ را حساب کنید.
\begin{solution}
ناحیه: مثلث پایین خط $x+y=1$ در مربع. برش عمودی: $y$ از $0$ تا $1-x$:
\[
\int_0^1\!\!\int_0^{1-x}(x+y)\dd y\dd x=\int_0^1\Big[x(1-x)+\frac{(1-x)^2}{2}\Big]\dd x=\int_0^1\frac{(1-x)(1+x)}{2}\dd x=\frac12\int_0^1(1-x^2)\dd x=\frac13 .
\]
(در قدم دوم $(1-x)$ را فاکتور گرفتیم: $(1-x)\big(x+\frac{1-x}{2}\big)=(1-x)\frac{1+x}{2}$.)
\end{solution}
\end{example}

\section{تغییر ترتیب انتگرال‌گیری}
یک ناحیه را می‌توان هم با برش عمودی و هم با برش افقی توصیف کرد. گاهی انتگرال در یک ترتیب ناممکن یا دشوار و در ترتیب دیگر ساده است.

\begin{example}[انتگرالی که فقط با تغییر ترتیب حل می‌شود]
$I=\int_0^1\int_x^1e^{-y^2}\dd y\dd x$ را حساب کنید.
\begin{solution}
انتگرال درونی $\int e^{-y^2}\dd y$ تابع اولیه مقدماتی ندارد؛ پس ترتیب را عوض می‌کنیم. \textbf{ناحیه:} $0\le x\le1$ و $x\le y\le1$، یعنی مثلث $0\le x\le y\le1$. \textbf{برش افقی:} به ازای $y$ ثابت در $[0,1]$، $x$ از $0$ تا $y$:
\[
I=\int_0^1\!\!\int_0^ye^{-y^2}\dd x\dd y=\int_0^1ye^{-y^2}\dd y=\Big[-\frac12e^{-y^2}\Big]_0^1=\frac{1-e^{-1}}{2}\approx0.316 .
\]
انتگرال درونی جدید ساده بود، چون انتگرالده به $x$ وابسته نیست و فقط طول بازه ($y$) ضرب می‌شود؛ این طول دقیقاً همان ضریب $y$ لازم برای جانشانی است.
\end{solution}
\end{example}

\begin{center}
\begin{tikzpicture}
\end{tikzpicture}\\
{\small یک ناحیه، دو توصیف: برش عمودی ($x$ ثابت، $y$ از $x$ تا ۱) و برش افقی ($y$ ثابت، $x$ از ۰ تا $y$). هر دو همان مثلث را می‌پوشانند.}
\end{center}

\begin{example}[اثبات دوم فرمول دم]
با تغییر ترتیب انتگرال، نشان دهید برای $X\ge0$: $\E X=\int_0^\infty\Prob(X>t)\dd t$.
\begin{solution}
ایده: $x=\int_0^x\dd t$. پس
\[
\E X=\int_0^\infty xf(x)\dd x=\int_0^\infty\Big(\int_0^x\dd t\Big)f(x)\dd x=\iint_{0\le t\le x}f(x)\dd t\dd x .
\]
ناحیه $0\le t\le x<\infty$ را با برش در جهت دیگر توصیف می‌کنیم: به ازای $t$ ثابت، $x$ از $t$ تا $\infty$:
\[
\E X=\int_0^\infty\Big(\int_t^\infty f(x)\dd x\Big)\dd t=\int_0^\infty\Prob(X>t)\dd t .
\]
(تغییر ترتیب برای توابع نامنفی همواره مجاز است؛ قضیه تونلی.) نسخه گسسته: برای $X$ با مقادیر صحیح نامنفی، $\E X=\sum_{k=1}^\infty\Prob(X\ge k)$، که با همین ترفند روی مجموع‌ها ثابت می‌شود.
\end{solution}
\end{example}

\section{مثال‌های کلاسیک احتمال}

\begin{example}[مسابقه دو ساعت نمایی]
$X\sim\mathrm{Exp}(\lambda)$ و $Y\sim\mathrm{Exp}(\mu)$ مستقل‌اند. $\Prob(X<Y)$ را بیابید.
\begin{solution}
چگالی مشترک (به دلیل استقلال) حاصل‌ضرب است: $\lambda e^{-\lambda x}\mu e^{-\mu y}$ برای $x,y\ge0$. ناحیه $x<y$: به ازای $x$ ثابت، $y$ از $x$ تا $\infty$:
\[
\Prob(X<Y)=\int_0^\infty\lambda e^{-\lambda x}\Big(\int_x^\infty\mu e^{-\mu y}\dd y\Big)\dd x=\int_0^\infty\lambda e^{-\lambda x}e^{-\mu x}\dd x=\frac{\lambda}{\lambda+\mu}.
\]
شهود: اگر دو رویداد با نرخ‌های $\lambda$ و $\mu$ «در رقابت» باشند، احتمال اینکه اولی زودتر رخ دهد متناسب با نرخ آن است. انتگرال درونی همان $\Prob(Y>x)=e^{-\mu x}$ است.
\end{solution}
\end{example}

\begin{example}[مسئله قرار ملاقات]
دو نفر مستقل از هم و به طور یکنواخت در یک ساعت (بازه $[0,1]$) به محل قرار می‌رسند و هر کدام فقط ۱۵ دقیقه منتظر می‌ماند. احتمال ملاقات چقدر است؟
\begin{solution}
$X,Y\sim\mathcal U(0,1)$ مستقل، با چگالی مشترک $1$ روی مربع واحد؛ پس احتمال هر ناحیه برابر \textbf{مساحت} آن است. ملاقات یعنی $|X-Y|\le\frac14$: نواری حول قطر. متمم آن دو مثلث قائم‌الزاویه در گوشه‌هاست ($y>x+\frac14$ و $x>y+\frac14$)، هر کدام با ساق‌های $\frac34$:
\[
\Prob\Big(|X-Y|\le\frac14\Big)=1-2\cdot\frac12\Big(\frac34\Big)^2=1-\frac{9}{16}=\frac{7}{16}.
\]
\end{solution}
\end{example}

\begin{center}
\begin{tikzpicture}
\end{tikzpicture}\\
{\small مسئله قرار ملاقات: نوار $|x-y|\le\frac14$ حول قطر مربع واحد. مساحت نوار برابر یک منهای مساحت دو مثلث گوشه است.}
\end{center}

\begin{example}[جمع دو متغیر یکنواخت مستقل]
$X,Y\sim\mathcal U(0,1)$ مستقل‌اند. تابع توزیع و چگالی $Z=X+Y$ را بیابید.
\begin{solution}
$F_Z(z)=\Prob(X+Y\le z)$ = مساحت بخشی از مربع واحد که زیر خط $x+y=z$ است. شکل ناحیه با $z$ عوض می‌شود، پس دو حالت داریم:
\begin{itemize}
\item $0\le z\le1$: مثلثی با ساق‌های $z$؛ پس $F_Z(z)=\frac{z^2}{2}$.
\item $1\le z\le2$: کل مربع منهای مثلث بالایی با ساق‌های $2-z$؛ پس $F_Z(z)=1-\frac{(2-z)^2}{2}$.
\end{itemize}
با مشتق‌گیری: $f_Z(z)=z$ برای $0\le z\le1$ و $f_Z(z)=2-z$ برای $1\le z\le2$؛ یعنی \textbf{توزیع مثلثی} (همان نتیجه کانولوشن جلسه ۱۶).
\end{solution}
\end{example}

\begin{center}
\begin{tikzpicture}
\end{tikzpicture}\\
{\small ناحیه $x+y\le z$ داخل مربع واحد برای $z=0.6$ (یک مثلث) و $z=1.4$ (مربع منهای یک مثلث). تغییر شکل ناحیه دلیل چندضابطه‌ای بودن جواب است.}
\end{center}

\begin{example}[جمع دو نمایی مستقل]
$X,Y\sim\mathrm{Exp}(\lambda)$ مستقل‌اند. چگالی $Z=X+Y$ را بیابید.
\begin{solution}
برای $z\ge0$، ناحیه $x,y\ge0$ و $x+y\le z$: به ازای $x\in[0,z]$، $y$ از $0$ تا $z-x$:
\[
F_Z(z)=\int_0^z\lambda e^{-\lambda x}\Big(\int_0^{z-x}\lambda e^{-\lambda y}\dd y\Big)\dd x=\int_0^z\lambda e^{-\lambda x}\big(1-e^{-\lambda(z-x)}\big)\dd x=\big(1-e^{-\lambda z}\big)-\lambda ze^{-\lambda z}.
\]
(در جمله دوم، $e^{-\lambda x}e^{-\lambda(z-x)}=e^{-\lambda z}$ ثابت است و انتگرالش $\lambda ze^{-\lambda z}$ می‌شود.) با مشتق‌گیری:
\[
f_Z(z)=\lambda e^{-\lambda z}-\lambda e^{-\lambda z}+\lambda^2ze^{-\lambda z}=\lambda^2ze^{-\lambda z},\qquad z\ge0,
\]
که چگالی گامای $(2,\lambda)$ است.
\end{solution}
\end{example}

\begin{example}[میانگین بیشینه دو یکنواخت با تقارن]
برای $X,Y\sim\mathcal U(0,1)$ مستقل، $\E\max(X,Y)$ را حساب کنید.
\begin{solution}
$\max(x,y)$ روی دو نیمه مربع ضابطه متفاوت دارد؛ پس انتگرال را روی قطر می‌شکنیم. به دلیل تقارن، دو نیمه سهم برابر دارند:
\[
\E\max(X,Y)=\iint_{y<x}x\dd y\dd x+\iint_{x<y}y\dd x\dd y=2\int_0^1\!\!\int_0^xx\dd y\dd x=2\int_0^1x^2\dd x=\frac23 .
\]
(بررسی: با فرمول $\frac{n}{n+1}$ برای $n=2$ یکی است.)
\end{solution}
\end{example}

\section{مختصات قطبی}
وقتی ناحیه یا انتگرالده به $x^2+y^2$ (فاصله تا مبدأ) وابسته است، مختصات قطبی کار را بسیار ساده می‌کند:
\[
x=r\cos\theta,\qquad y=r\sin\theta,\qquad r\ge0,\ \ 0\le\theta<2\pi,\qquad x^2+y^2=r^2 .
\]
\begin{theorem}[عنصر مساحت قطبی]
\[
\dd x\dd y=r\dd r\dd\theta .
\]
\end{theorem}
\textbf{شهود:} یک «مستطیل قطبی» کوچک بین شعاع‌های $r$ و $r+\dd r$ و زاویه‌های $\theta$ و $\theta+\dd\theta$ تقریباً یک مستطیل با ضلع شعاعی $\dd r$ و ضلع کمانی $r\dd\theta$ است (طول کمان = شعاع × زاویه). پس مساحتش $r\dd r\dd\theta$ است. ضریب $r$ می‌گوید: سلول‌های دورتر از مبدأ، با همان $\dd r$ و $\dd\theta$، بزرگ‌ترند.

\begin{center}
\begin{tikzpicture}
\end{tikzpicture}\\
{\small عنصر مساحت در مختصات قطبی: ضلع شعاعی $\dd r$ و ضلع کمانی $r\,\dd\theta$؛ پس مساحت سلول $r\,\dd r\,\dd\theta$ است.}
\end{center}

\begin{theorem}[اثبات انتگرال گاوسی]
$I=\int_{-\infty}^{\infty}e^{-x^2/2}\dd x=\sqrt{2\pi}$.
\end{theorem}
\begin{proof}
ترفند: به جای $I$، مربع آن را حساب می‌کنیم و آن را یک انتگرال دوگانه می‌نویسیم (فوبینی، حالت جداپذیر):
\[
I^2=\int_{-\infty}^{\infty}e^{-x^2/2}\dd x\int_{-\infty}^{\infty}e^{-y^2/2}\dd y=\iint_{\R^2}e^{-(x^2+y^2)/2}\dd x\dd y .
\]
انتگرالده فقط به $r^2=x^2+y^2$ وابسته است؛ به قطبی می‌رویم (کل صفحه: $r\in[0,\infty)$ و $\theta\in[0,2\pi)$):
\[
I^2=\int_0^{2\pi}\!\!\int_0^\infty e^{-r^2/2}\,r\dd r\dd\theta=2\pi\cdot\Big[-e^{-r^2/2}\Big]_0^\infty=2\pi\cdot1 .
\]
ضریب $r$ از عنصر مساحت، دقیقاً همان چیزی است که جانشانی $u=\frac{r^2}{2}$ لازم دارد! چون $I>0$، $I=\sqrt{2\pi}$.
\end{proof}

\begin{example}[فاصله یک نقطه گاوسی از مبدأ]
$X,Y\sim\mathcal N(0,1)$ مستقل‌اند و $R=\sqrt{X^2+Y^2}$. توزیع $R$ و $R^2$ را بیابید.
\begin{solution}
چگالی مشترک $\frac{1}{2\pi}e^{-(x^2+y^2)/2}$ است. ناحیه $R\le\rho$ یک دایره است که در قطبی یعنی $0\le r\le\rho$ و هر $\theta$:
\[
\Prob(R\le\rho)=\int_0^{2\pi}\!\!\int_0^\rho\frac{1}{2\pi}e^{-r^2/2}\,r\dd r\dd\theta=1-e^{-\rho^2/2},\qquad\rho\ge0 .
\]
با مشتق‌گیری: $f_R(\rho)=\rho e^{-\rho^2/2}$ (توزیع \textbf{رایلی}). همچنین $\Prob(R^2\le s)=\Prob(R\le\sqrt s)=1-e^{-s/2}$، یعنی $R^2=X^2+Y^2\sim\mathrm{Exp}(\frac12)$ که همان $\chi^2(2)$ است (جلسه ۱۷).
\end{solution}
\end{example}

\begin{example}[تقارن چرخشی]
برای $X,Y\sim\mathcal N(0,1)$ مستقل، $\Prob(0<Y<X)$ را بدون انتگرال‌گیری مستقیم بیابید.
\begin{solution}
چگالی فقط به $r$ وابسته است، پس زاویه $\theta$ نقطه $(X,Y)$ روی $[0,2\pi)$ \textbf{یکنواخت} است (از انتگرال قطبی: سهم هر بازه زاویه‌ای متناسب با طول آن است). ناحیه $0<y<x$ یعنی $0<\theta<\frac\pi4$:
\[
\Prob(0<Y<X)=\frac{\pi/4}{2\pi}=\frac18 .
\]
\end{solution}
\end{example}

\section{تغییر متغیر کلی و ژاکوبین}

\begin{theorem}[تغییر متغیر در انتگرال دوگانه]
فرض کنید نگاشت $(u,v)\mapsto\big(x(u,v),y(u,v)\big)$ یک‌به‌یک و مشتق‌پذیر باشد و ناحیه $B$ در صفحه $(u,v)$ را به ناحیه $A$ در صفحه $(x,y)$ ببرد. آنگاه
\[
\iint_Af(x,y)\dd x\dd y=\iint_Bf\big(x(u,v),y(u,v)\big)\,|J|\dd u\dd v,\qquad
J=\det\begin{bmatrix}\dfrac{\partial x}{\partial u}&\dfrac{\partial x}{\partial v}\\[3mm]\dfrac{\partial y}{\partial u}&\dfrac{\partial y}{\partial v}\end{bmatrix}.
\]
\end{theorem}
\textbf{یادآوری:} $\det\begin{bmatrix}a&b\\c&d\end{bmatrix}=ad-bc$.

\textbf{شهود:} $|J|$ \textbf{ضریب تغییر مساحت} نگاشت است: یک مستطیل کوچک $\dd u\times\dd v$ در صفحه $(u,v)$ به یک متوازی‌الاضلاع کوچک در صفحه $(x,y)$ نگاشته می‌شود که مساحتش $|J|\dd u\dd v$ است. (دترمینان، مساحت متوازی‌الاضلاعی است که دو بردار ستونی ماتریس می‌سازند.) در حالت یک‌بعدی این همان $|\frac{\dd x}{\dd u}|$ در جانشانی است. برای قطبی: $J=\det\begin{bmatrix}\cos\theta&-r\sin\theta\\\sin\theta&r\cos\theta\end{bmatrix}=r\cos^2\theta+r\sin^2\theta=r$؛ همان عنصر مساحت بخش قبل.

\subsection{کاربرد اصلی: چگالی تبدیل دو متغیر تصادفی}
اگر $(U,V)=g(X,Y)$ یک‌به‌یک باشد و وارون آن $x=x(u,v)$، $y=y(u,v)$ باشد، آنگاه (با همان استدلال حفظ احتمال هر ناحیه کوچک):
\[
\boxed{\ f_{U,V}(u,v)=f_{X,Y}\big(x(u,v),y(u,v)\big)\,\big|J\big|,\qquad J=\frac{\partial(x,y)}{\partial(u,v)}\ }
\]
معادلاً $|J|=\Big|\det\frac{\partial(u,v)}{\partial(x,y)}\Big|^{-1}$، یعنی می‌توان ژاکوبین نگاشت مستقیم را حساب کرد و وارونش را گرفت (گاهی آسان‌تر است). این تعمیم دوبعدی فرمول $f_Y(y)=f_X(x)\big|\frac{\dd x}{\dd y}\big|$ جلسه ۱۲ است.

\begin{example}[جمع و تفاضل دو گاوسی مستقل]
$X,Y\sim\mathcal N(0,1)$ مستقل‌اند، $U=X+Y$ و $V=X-Y$. چگالی مشترک $(U,V)$ را بیابید.
\begin{solution}
\textbf{وارون:} $x=\frac{u+v}{2}$ و $y=\frac{u-v}{2}$.
\textbf{ژاکوبین:}
\[
J=\det\begin{bmatrix}\frac12&\frac12\\[1mm]\frac12&-\frac12\end{bmatrix}=-\frac14-\frac14=-\frac12,\qquad|J|=\frac12 .
\]
\textbf{جایگذاری:} $x^2+y^2=\frac{(u+v)^2+(u-v)^2}{4}=\frac{u^2+v^2}{2}$؛ پس
\[
f_{U,V}(u,v)=\frac{1}{2\pi}e^{-\frac{u^2+v^2}{4}}\cdot\frac12=\Big(\frac{1}{\sqrt{4\pi}}e^{-u^2/4}\Big)\Big(\frac{1}{\sqrt{4\pi}}e^{-v^2/4}\Big).
\]
چگالی جداپذیر شد: $U$ و $V$ \textbf{مستقل} و هر دو $\mathcal N(0,2)$ هستند. (برای گاوسی‌ها ناهمبستگی معادل استقلال است: $\Cov(X+Y,X-Y)=\Var X-\Var Y=0$؛ جلسه ۱۴.)
\end{solution}
\end{example}

\begin{example}[جمع و نسبت دو نمایی]
$X,Y\sim\mathrm{Exp}(1)$ مستقل‌اند، $U=X+Y$ و $V=\frac{X}{X+Y}$. چگالی مشترک $(U,V)$ را بیابید.
\begin{solution}
\textbf{وارون:} $x=uv$ و $y=u-x=u(1-v)$، با $u>0$ و $0<v<1$.
\textbf{ژاکوبین:}
\[
J=\det\begin{bmatrix}v&u\\1-v&-u\end{bmatrix}=-uv-u(1-v)=-u,\qquad|J|=u .
\]
\textbf{جایگذاری:} $f_{X,Y}(x,y)=e^{-x-y}=e^{-u}$:
\[
f_{U,V}(u,v)=e^{-u}\cdot u=\underbrace{ue^{-u}}_{u>0}\cdot\underbrace{1}_{0<v<1}.
\]
پس $U\sim\mathrm{Gamma}(2,1)$ (همان نتیجه جمع دو نمایی)، $V\sim\mathcal U(0,1)$ و این دو مستقل‌اند! یعنی «سهم $X$ از کل زمان» مستقل از کل زمان و یکنواخت است.
\end{solution}
\end{example}

\begin{theorem}[رابطه بتا و گاما]
برای $a,b>0$: $B(a,b)=\dfrac{\Gamma(a)\Gamma(b)}{\Gamma(a+b)}$.
\end{theorem}
\begin{proof}
حاصل‌ضرب دو گاما را یک انتگرال دوگانه می‌نویسیم و همان تغییر متغیر مثال قبل ($x=uv$، $y=u(1-v)$، $|J|=u$) را اعمال می‌کنیم:
\[
\Gamma(a)\Gamma(b)=\int_0^\infty\!\!\int_0^\infty x^{a-1}y^{b-1}e^{-(x+y)}\dd x\dd y
=\int_0^\infty\!\!\int_0^1(uv)^{a-1}\big(u(1-v)\big)^{b-1}e^{-u}\,u\dd v\dd u .
\]
توان‌های $u$: $(a-1)+(b-1)+1=a+b-1$. پس انتگرال جداپذیر می‌شود:
\[
\Gamma(a)\Gamma(b)=\int_0^\infty u^{a+b-1}e^{-u}\dd u\cdot\int_0^1v^{a-1}(1-v)^{b-1}\dd v=\Gamma(a+b)\,B(a,b).
\]
\end{proof}

\section{انتگرال سه‌گانه (اشاره)}
همه ایده‌ها به سه بعد (و بیشتر) تعمیم می‌یابند: $\iiint f\dd x\dd y\dd z$ به صورت سه انتگرال تکراری حساب می‌شود و ژاکوبین یک دترمینان $3\times3$ است.

\begin{example}[حجم سادک]
$X,Y,Z\sim\mathcal U(0,1)$ مستقل‌اند. (الف) $\Prob(X+Y+Z\le1)$ و (ب) $\Prob(X<Y<Z)$ را حساب کنید.
\begin{solution}
چگالی مشترک روی مکعب واحد برابر ۱ است؛ پس احتمال‌ها \textbf{حجم}اند.

(الف) به ازای $x$ و $y$ ثابت، $z$ از $0$ تا $1-x-y$؛ به ازای $x$ ثابت، $y$ از $0$ تا $1-x$:
\[
\int_0^1\!\!\int_0^{1-x}(1-x-y)\dd y\dd x=\int_0^1\frac{(1-x)^2}{2}\dd x=\frac16 .
\]
(انتگرال درونی: مساحت مثلثی با ساق‌های $1-x$.) به طور کلی برای $n$ متغیر، $\Prob(X_1+\cdots+X_n\le1)=\frac{1}{n!}$.

(ب) $\int_0^1\int_0^z\int_0^y\dd x\dd y\dd z=\int_0^1\int_0^zy\dd y\dd z=\int_0^1\frac{z^2}{2}\dd z=\frac16$. با تقارن هم واضح است: هر یک از $3!=6$ ترتیب ممکن هم‌احتمال است.
\end{solution}
\end{example}

\begin{summarybox}[جمع‌بندی: نقشه راه انتگرال دوگانه در احتمال]
\begin{enumerate}
\item همه شرط‌ها (تکیه‌گاه چگالی + رویداد) را روی صفحه رسم و اشتراک را سایه بزنید.
\item اگر ناحیه دایره‌ای یا انتگرالده تابعی از $x^2+y^2$ است: قطبی با $\dd x\dd y=r\dd r\dd\theta$.
\item وگرنه یک ترتیب انتخاب کنید؛ حدود درونی تابعی از متغیر بیرونی، حدود بیرونی عدد ثابت.
\item اگر انتگرال درونی سخت است، ترتیب را عوض کنید. اگر شکل ناحیه با پارامتر ($z$) عوض می‌شود، حالت‌بندی کنید.
\item برای چگالی یک تبدیل یک‌به‌یک: وارون را بنویسید، $|J|$ را حساب کنید و $f_{U,V}=f_{X,Y}\cdot|J|$؛ حتماً ناحیه جدید $(u,v)$ را هم مشخص کنید.
\item چگالی جداپذیر روی ناحیه مستطیلی (حاصل‌ضربی) یعنی استقلال.
\end{enumerate}
\end{summarybox}

\section{تمرین‌های تکمیلی}

\begin{exercise}[نمایی‌های مستقل]
$f(x,y)=ce^{-x-2y}$ برای $x,y\ge0$. (الف) $c$ را بیابید. (ب) آیا $X,Y$ مستقل‌اند؟ (ج) $\Prob(X<Y)$ را حساب کنید.
\end{exercise}
\begin{solution}
(الف) جداپذیر روی ناحیه حاصل‌ضربی: $\int_0^\infty e^{-x}\dd x\int_0^\infty e^{-2y}\dd y=1\cdot\frac12$؛ پس $c=2$.
(ب) بله: $f=\big(e^{-x}\big)\big(2e^{-2y}\big)$، یعنی $X\sim\mathrm{Exp}(1)$ و $Y\sim\mathrm{Exp}(2)$ مستقل.
(ج) با فرمول مسابقه: $\frac{1}{1+2}=\frac13$. مستقیم: $\int_0^\infty e^{-x}\int_x^\infty2e^{-2y}\dd y\dd x=\int_0^\infty e^{-3x}\dd x=\frac13$.
\end{solution}

\begin{exercise}[تکیه‌گاه غیرمستطیلی]
$f(x,y)=8xy$ برای $0<y<x<1$. (الف) نرمال بودن را بررسی و حاشیه‌ای‌ها را حساب کنید. (ب) با اینکه $8xy$ ظاهراً جداپذیر است، آیا $X,Y$ مستقل‌اند؟
\end{exercise}
\begin{solution}
(الف) $\int_0^1\int_0^x8xy\dd y\dd x=\int_0^14x^3\dd x=1$. $f_X(x)=\int_0^x8xy\dd y=4x^3$ و $f_Y(y)=\int_y^18xy\dd x=4y(1-y^2)$.
(ب) خیر: $f_X(x)f_Y(y)=16x^3y(1-y^2)\neq8xy$. دلیل عمیق‌تر: تکیه‌گاه مثلث است نه مستطیل؛ دانستن $X=0.2$ می‌گوید $Y<0.2$، پس $Y$ به $X$ وابسته است. چگالی در واقع $8xy\,\ind_{y<x}$ است و نشانگر $\ind_{y<x}$ جداپذیر نیست.
\end{solution}

\begin{exercise}[تغییر ترتیب]
$\int_0^1\int_y^1\frac{\sin x}{x}\dd x\dd y$ را حساب کنید.
\end{exercise}
\begin{solution}
$\frac{\sin x}{x}$ تابع اولیه مقدماتی ندارد. ناحیه $0\le y\le x\le1$؛ با برش عمودی ($x$ ثابت، $y$ از $0$ تا $x$):
\[
\int_0^1\!\!\int_0^x\frac{\sin x}{x}\dd y\dd x=\int_0^1\frac{\sin x}{x}\cdot x\dd x=\int_0^1\sin x\dd x=1-\cos1\approx0.460 .
\]
\end{solution}

\begin{exercise}[ملاقات با ۱۰ دقیقه انتظار]
در مسئله قرار ملاقات، اگر هر نفر ۱۰ دقیقه منتظر بماند، احتمال ملاقات چقدر است؟ به طور کلی با زمان انتظار $w$ (کسری از ساعت)؟
\end{exercise}
\begin{solution}
$1-(1-w)^2$. برای $w=\frac16$: $1-\frac{25}{36}=\frac{11}{36}\approx0.306$.
\end{solution}

\begin{exercise}[نقطه یکنواخت در دایره]
نقطه $(X,Y)$ به طور یکنواخت در دایره واحد انتخاب می‌شود (چگالی $\frac1\pi$). توزیع فاصله $R$ از مرکز و $\E R$ را بیابید.
\end{exercise}
\begin{solution}
$\Prob(R\le r)=\frac{\pi r^2}{\pi}=r^2$ (نسبت مساحت‌ها)، پس $f_R(r)=2r$ برای $0\le r\le1$ و $\E R=\int_0^12r^2\dd r=\frac23$. توجه: $R$ یکنواخت نیست؛ نقاط دور از مرکز بیشترند، چون عنصر مساحت $r\dd r\dd\theta$ با $r$ رشد می‌کند.
\end{solution}

\begin{exercise}[نسبت دو گاوسی]
$X,Y\sim\mathcal N(0,1)$ مستقل‌اند. با استفاده از یکنواخت بودن زاویه قطبی، نشان دهید $V=\frac YX$ توزیع کوشی دارد.
\end{exercise}
\begin{solution}
$V=\frac{R\sin\Theta}{R\cos\Theta}=\tan\Theta$ با $\Theta\sim\mathcal U[0,2\pi)$. $\tan$ با دوره $\pi$ تکرار می‌شود، پس می‌توان $\Theta$ را یکنواخت روی $(-\frac\pi2,\frac\pi2)$ گرفت. آنگاه
\[
\Prob(V\le v)=\Prob(\Theta\le\arctan v)=\frac{\arctan v+\frac\pi2}{\pi}=\frac12+\frac{\arctan v}{\pi},
\]
که همان تابع توزیع کوشی (ریاضی ۱) است.
\end{solution}

\begin{exercise}[جمع و تفاضل دو یکنواخت]
$X,Y\sim\mathcal U(0,1)$ مستقل، $U=X+Y$ و $V=X-Y$. چگالی مشترک $(U,V)$ و ناحیه آن را بیابید و از آن چگالی $U$ را به دست آورید.
\end{exercise}
\begin{solution}
مانند مثال گاوسی، $|J|=\frac12$؛ پس $f_{U,V}=1\cdot\frac12=\frac12$ روی تصویر مربع. شرط‌های $0<x,y<1$ با $x=\frac{u+v}{2}$ و $y=\frac{u-v}{2}$ یعنی $|v|<u$ و $|v|<2-u$: یک \textbf{مربع چرخیده} (لوزی) با رئوس $(0,0)$، $(1,1)$، $(2,0)$ و $(1,-1)$. حاشیه‌ای: برای $0<u<1$، $v$ از $-u$ تا $u$، پس $f_U(u)=\frac12\cdot2u=u$؛ برای $1<u<2$، $v$ از $-(2-u)$ تا $2-u$، پس $f_U(u)=2-u$. همان توزیع مثلثی.
\end{solution}
